\section{Additive and spectral extensions}
\label{sec:motivic}

The blowup formulas also define an additive invariant on the Grothendieck
group of varieties.  We first retain the squared exponent gaps instead of
only counting them.  The canonical modified lattice is essential here:
arbitrary separate integer shifts of the two exponents can change their
squared difference.

\begin{lemma}[A fixed target for the residue spectrum]
\label{lem:constant-residue-spectrum}
\coverage{absent}%
\uses{prop:rank2-rigidity,lem:faithful-center-base-change,prop:qdm-operation-ledgers}%
For every counted generic block, \(\delta^\sharp\) belongs to
\[
 \mathcal D=\C\setminus\{n^2:n\in\Z\}.
\]
Its value is preserved by the regular comparisons used in the blowup and
projective-bundle formulas.
\end{lemma}

\begin{proof}
The finite jet defining the residue is over a finite algebraic extension
of the reduced generic coefficient field \(K_Y\): spectral projectors
and the finitely many Sylvester equations needed for this jet require
only algebraic scalar extension.  Proposition~\ref{prop:rank2-rigidity}
makes its discriminant constant in every even bulk direction.  On the
reduced source, a divisor direction acts by
\(X_d\mapsto(D\cdot d)X_d\), while the other bulk directions are ordinary
partial derivatives.  Thus the discriminant is killed by all logarithmic
derivations in the numerical curve and remaining bulk variables.

The common constants of \(K_Y\) under these derivations are \(\C\).
Here is a coefficient argument that also applies to the completed ring.
Write a constant fraction as \(f/g\).  Simultaneous formal rescaling of
the monomials gives \(f(e^t x)g(x)=f(x)g(e^t x)\), by uniqueness of formal
Taylor expansion.  At each bounded ample-plus-bulk degree only finitely
many monomials occur.  Their distinct exponential characters are linearly
independent, by the Vandermonde argument in
Lemma~\ref{lem:faithful-center-base-change}.  Comparing the coefficient
of each monomial therefore gives
\(f_a g_b=g_a f_b\) for every pair of indices \(a,b\).
Choose \(b\) with \(g_b\ne0\); then \(f=(f_b/g_b)g\).
The same constant field holds in an algebraic extension: differentiating
the monic minimal polynomial of an algebraic constant gives a polynomial
of smaller degree vanishing at it, so every coefficient of the minimal
polynomial is constant.  Since \(\C\) is algebraically closed, the
algebraic constant lies in \(\C\).  The exponent criterion excludes
exactly the integer squares.

A regular comparison maps the centered leading operator and its image
line to their counterparts.  It therefore maps \(E^\sharp\) isomorphically
to the corresponding modified lattice.  On reduction modulo \(z\), the
residues are conjugate.  A common scalar shift of the residue also leaves
\((\tr R)^2-4\det R\) unchanged.  These are the transformations supplied
by the comparisons; no separate change of the two exponent representatives
is needed.  The faithful coefficient maps fix \(\C\), so they preserve
the same numerical value in \(\mathcal D\).
\uses{prop:rank2-rigidity,lem:faithful-center-base-change,prop:qdm-operation-ledgers}%
\end{proof}

Let \(M_{\mathrm{disc}}=\bigoplus_{\delta\in\mathcal D}\mathbf N e_\delta\), with group completion
\(G_{\mathrm{disc}}=\bigoplus_{\delta\in\mathcal D}\Z e_\delta\).  For a connected smooth projective
variety put
\[
 \mathfrak S(Y)=\sum_{E\text{ counted}}e_{\delta^\sharp(E)}\in M_{\mathrm{disc}}.
\]
Extend by sums over connected components and put
\(\mathfrak S(\varnothing)=0\).  The augmentation
\(\epsilon:G_{\mathrm{disc}}\to\Z\), \(e_\delta\mapsto1\), recovers
\(I_{\mathrm{exp}}\).

\begin{proposition}[Intrinsic operation formulas]
\label{prop:intrinsic-spectrum-formulas}
\coverage{absent}%
\uses{lem:constant-residue-spectrum,lem:faithful-center-base-change,prop:qdm-operation-ledgers}%
Let \(Y\) be smooth projective over \(\C\).  For a smooth closed center
\(Z\subset Y\) of codimension \(c\geq1\) and
a rank-\(r\) vector bundle \(V\) on \(Y\),
\[
 \mathfrak S(\Bl_ZY)=\mathfrak S(Y)+(c-1)\mathfrak S(Z),
 \qquad \mathfrak S(\PP_Y(V))=r\mathfrak S(Y).
\]
The same formulas hold for \(I_{\mathrm{exp}}\).  For centers of differing
component dimensions, apply the formula separately to each component.
\end{proposition}

\begin{proof}
Lemma~\ref{lem:faithful-center-base-change} identifies every center
occurrence with a faithful extension of its intrinsic module in every
dimension.  Lemma~\ref{lem:constant-residue-spectrum} preserves each
summand's spectrum.  Independent unit coordinates separate the summands
generically, as in Proposition~\ref{prop:qdm-operation-ledgers}.
Consequently their spectra add.  The projective-bundle comparison gives
\(r\) copies in the same way.  Augmentation proves the scalar formulas.
Disjoint centers may be blown up successively, giving the last assertion.
\uses{lem:constant-residue-spectrum,lem:faithful-center-base-change,prop:qdm-operation-ledgers}%
\end{proof}

\begin{theorem}[Additive Grothendieck-group invariants]
\label{thm:additive-motivic-marker}
\coverage{absent}%
\uses{prop:intrinsic-spectrum-formulas}%
\imports{Bittner:blowup-presentation}%
There are unique additive homomorphisms
\[
 \mathcal S:K_0(\operatorname{Var}_{\C})/(\mathbb L-1)\longrightarrow G_{\mathrm{disc}},
 \qquad
 \mathcal I=\epsilon\mathcal S:
 K_0(\operatorname{Var}_{\C})/(\mathbb L-1)\longrightarrow\Z
\]
agreeing with \(\mathfrak S\) and \(I_{\mathrm{exp}}\) on smooth projective
varieties.  Here \(\mathbb L=[\mathbf A^1]\),
\(\mathcal I(1)=\mathcal I(\mathbb L)=0\), and for a smooth cubic
threefold \(X\), \(\mathcal S([X])=e_{4/9}\) and \(\mathcal I([X])=1\).
\end{theorem}

\begin{proof}
Bittner's presentation of the abelian group uses smooth projective
generators, the empty class, and the blowup relation
\cite[Theorem~3.1]{Bittner}.  Since the exceptional divisor is
\(\PP_Z(N_{Z/Y})\), Proposition~\ref{prop:intrinsic-spectrum-formulas} gives
\[
 \mathfrak S(\Bl_ZY)-\mathfrak S(\PP_Z(N_{Z/Y}))
 =\mathfrak S(Y)-\mathfrak S(Z)
\]
in \(G_{\mathrm{disc}}\).  Thus the assignment extends uniquely and additively.
For a smooth projective \(Y\), the scissor relation and the bundle formula
give
\[
 [Y\times\PP^1]=[Y]+\mathbb L[Y],\qquad
 \mathcal S(\mathbb L[Y])=
 \mathfrak S(Y\times\PP^1)-\mathfrak S(Y)=\mathfrak S(Y).
\]
Smooth projective classes generate the whole group, so
\(\mathcal S((\mathbb L-1)a)=0\) for every \(a\).  This proves the
factorization through the indicated quotient, and augmentation gives
\(\mathcal I\).  Points have no counted block, while the cubic has the
single block of discriminant \(4/9\) computed in
Section~\ref{subsec:shared-cubic-packet}.
\uses{prop:intrinsic-spectrum-formulas,prop:atomic-lowdim}%
\imports{Bittner:blowup-presentation}%
\end{proof}

These are additive group invariants.  In particular, \(\mathcal I\) is
not multiplicative: its value at \(1\) is zero but its value at a cubic
is one.  The identity \(\mathcal I(\mathbb L a)=\mathcal I(a)\) does not make
\(\mathcal I\) a stable-birational invariant.

For \(n\geq1\), the relation \(H^{n+1}=q\) in small quantum cohomology
shows that Euler multiplication on \(\PP^n\) has distinct eigenvalues at
\(q\ne0\).  It remains semisimple at the generic big point.  Thus
\(\mathfrak S(\PP^n)=0\) for every \(n\), including \(n=0\).
For example, embed a cubic threefold \(X\subset\PP^4\subset\PP^5\).
Then \(\mathcal I([\Bl_X\PP^5])=1\), whereas
\(\mathcal I([\PP^5])=0\).  Thus birational invariance already fails
in dimension five.

\begin{proposition}[Equal Hodge diamonds, different invariant values]
\label{prop:hodge-marker-separation}
\coverage{absent}%
\uses{prop:intrinsic-spectrum-formulas,prop:atomic-lowdim}%
Let \(X\) be a smooth cubic threefold and \(p\in X\).  Let
\(C\subset\PP^1\times\PP^1\subset\PP^3\) be a smooth curve of
bidegree \((2,6)\).  Then \(\Bl_pX\) and \(\Bl_C\PP^3\) have the same
Hodge diamond, but their invariant values are respectively one and zero.
Consequently \(\mathcal I\) does not factor through the Hodge--Deligne
polynomial.
\end{proposition}

\begin{proof}
The linear system \(\OO(2,6)\) contains smooth members, and adjunction
gives genus \((2-1)(6-1)=5\).  For the cubic,
\(c(T_X)=(1+H)^5/(1+3H)\) gives Euler characteristic \(-6\).
Weak Lefschetz and \(K_X=\OO_X(-2)\) then give
\(h^{1,1}=h^{2,2}=1\), \(h^{2,1}=h^{1,2}=5\), and
\(h^{3,0}=h^{0,3}=0\).  The Hodge blowup formula shows that both displayed
blowups have, besides \(h^{0,0}=h^{3,3}=1\), precisely
\[
 h^{1,1}=h^{2,2}=2,\qquad h^{2,1}=h^{1,2}=5.
\]
Their Hodge--Deligne polynomials therefore agree.  The intrinsic blowup
formula adds only point or curve contributions, all zero, so their invariant
values equal those of \(X\) and \(\PP^3\), namely one and zero.
\uses{prop:intrinsic-spectrum-formulas,prop:atomic-lowdim}%
\end{proof}

\begin{corollary}[Contributions in any rationalizing factorization]
\label{cor:factorization-spectrum}
\coverage{absent}%
\uses{prop:intrinsic-spectrum-formulas,prop:atomic-lowdim}%
\imports{AKMW:weak-factorization}%
Let \(X\) be a smooth cubic threefold.  In any projective weak
factorization of a rationalization
\(X\times\PP^m\dashrightarrow\PP^{m+3}\), orient the arrows from the
product to projective space and take connected centers.  Then
\[
 \sum_{\mathrm{blowdowns}}(c_Z-1)\mathfrak S(Z)
 -\sum_{\mathrm{blowups}}(c_Z-1)\mathfrak S(Z)
 =(m+1)e_{4/9}.
\]
For \(m=2\), only threefold centers of codimension two can contribute;
their net spectrum is \(3e_{4/9}\), and their net contribution to \(I_{\mathrm{exp}}\) is three.
\end{corollary}

\begin{proof}
The initial spectrum is \((m+1)e_{4/9}\); projective space has no counted
blocks and hence spectrum zero.  Telescope the intrinsic blowup formulas.
In dimension five, every nontrivial center has dimension at most three,
and centers of dimension at most two contribute zero.  The remaining
centers have codimension two.  Applying \(\epsilon\) gives the scalar
identity.  This counts contributions with multiplicity, not distinct
centers, and does not identify a center as a cubic threefold.
\uses{prop:intrinsic-spectrum-formulas,prop:atomic-lowdim}%
\imports{AKMW:weak-factorization}%
\end{proof}

\paragraph{Limits of the construction.}
Smooth complex cubic threefolds form one connected deformation family.
Unmarked deformation-invariant Gromov--Witten data are consequently the
same throughout that family, whereas it contains both stably rational
cubics \cite{TschinkelZhangTorsors} and very general stably irrational
cubics \cite{EGFS}.  Increasing the collection of such quantum data
cannot by itself completely detect stable rationality within the family.

The rank-two pairing argument also has a specific boundary.  In rank
three, let \(E_{ij}\) denote the matrix units and put
\[
 N=E_{12}+E_{23},\qquad
 P=\begin{pmatrix}0&0&1\\0&1&0\\1&0&0\end{pmatrix},\qquad
 A_0=E_{21}-E_{32}.
\]
Then
\[
 N^TP=PN=E_{23}+E_{32},\qquad
 PA_0=E_{21}-E_{12}=-A_0^TP.
\]
Thus the constant pairing is horizontal for \(A(z)=N/z+A_0\), but
\(A_0e_1=e_2\notin\ker N\).  Pairing horizontality alone does not
preserve the kernel flag.  For a cyclic Jordan leading term and
\(S=\diag(1,z,\ldots,z^{r-1})\), regularity of the modified connection
requires \((A_k)_{ij}=0\) whenever \(i-j>k\), for \(k\geq0\).
Additional geometric filtration compatibility would therefore be needed
for that higher-rank modification.
